% Part: computability
% Chapter: computability-theory
% Section: coding-computations

\documentclass[../../../include/open-logic-section]{subfiles}

\begin{document}

\olfileid{cmp}{thy}{cod}
\olsection{حساباں نوں کوڈ دینا}

حساب دے ہر نمونے وچ ایہ کم کیتے جا سکدے نیں:
\begin{enumerate}
\item قابلِ حساب فنکشناں دیاں \emph{تعریفاں} نوں اک منظم طریقے
  نال بیان کرنا۔ مثلاً ٹیورنگ مشین دیاں تعییناں، بازگشتی تعریفاں
  یا کسے پروگرامنگ زبان دے پروگرام ایہ تعریفاں دے سکدے نیں۔
\item کسے تعریف نال دِتے فنکشن دے کسے دِتے انپٹ اُتے حساب دا
  پورا ریکارڈ بیان کرنا۔ مثلاً ٹیورنگ مشین دا حساب اوہناں
  ترتیب وار حالتاں دے سلسلے نال بیان ہو سکدا اے جیہڑیاں حساب
  دے ہر قدم اُتے ہوندیاں نیں: مشین دی حالت تے ٹیپ دا مواد۔
\item ایہ پرکھنا کہ حساب دا دِتا ہویا ممکنہ ریکارڈ واقعی اوہ
  ریکارڈ اے جیہدے مطابق دِتی تعریف والا قابلِ حساب فنکشن دِتے
  انپٹ اُتے حسابیا جاندا اے؛ ایس انپٹ اُتے فنکشن تعریف شدہ
  ہوندا اے، یعنی حساب رُک جاندا اے۔
\item حساب دے پورے ریکارڈ دے ایہوجیہے بیان توں دِتے انپٹ لئی
  فنکشن دی قدر کڈھنا۔ مثلاً ٹیورنگ مشین دے رُکّن والے حساب
  دے بالکل آخری قدم اُتے ٹیپ دا مواد ایہ قدر ہوندا اے۔
\end{enumerate}

کوڈنگ راہیں قابلِ حساب فنکشن دے ہر بیان نوں اک عددی
\emph{اشاریہ (index)} دِتا جا سکدا اے، تے ایہ ایداں کیتا
جا سکدا اے کہ اشاریے توں ہدایتاں اک قابلِ حساب طریقے نال
واپس کڈھیاں جا سکّن۔ ایس طرح حساب دے پورے ریکارڈ نوں وی
اکّو عدد دے کوڈ وچ رکھیا جا سکدا اے۔ ایس توں بنن والا حسابی
تعلق ``$s$، اشاریہ~$e$ والے فنکشن دے انپٹ~$x$ اُتے حساب
دے ریکارڈ دا کوڈ اے'' تے فنکشن ``کوڈ~$s$ والے حسابی سلسلے
دا نتیجہ'' دونوں قابلِ حساب ہوندے نیں؛ بلکہ اوہ بنیادی
بازگشتی نیں۔

ایہ بنیادی حقیقت بڑی طاقتور اے، تے ایس دی مدد نال اسی
قابلِ حسابیت بارے کئی حیران کن تے اہم نتیجے ثابت کر سکدے
آں، بھاویں حساب دا جیہڑا وی نمونہ چنیا گیا ہووے۔

\end{document}
